\section{评估是在为一句关于未见数据的断言背书}
\label{sec:17-Synthesis/01-System-Lifecycle/05}

新的排序策略在离线评测上把 NDCG 提高了 \(2.1\%\)。上线做 A/B，跑满两周，\(p=0.12\)。产品经理说再跑一周看看；第三周末 \(p=0.043\)，会上决定全量。一个月后复盘，大盘指标与上线前没有可辨认的差别。

这个过程里每一步都符合流程，结论却是错的。问题出在``再跑一周看看''这个动作上：每天盯着 \(p\) 值、一旦跨过阈值就停止实验，这套程序的实际第一类错误率远高于名义上的 \(5\%\)。原本的承诺是``在原假设为真时，这套判定规则出错的概率不超过 \(5\%\)''，而边看边停已经不是原来那套规则了。

评估这一步的产出看起来只是几个数，实际上它是一句承诺：\textbf{这个系统在它没见过的数据上会怎样}。承诺是关于未来的，而手上只有过去的一次抽样，所以这句话能不能算数，取决于产生这几个数的\emph{程序}是否真的模拟了``未见''。

\subsection{置信度属于程序，不属于这次的参数}

先把最容易说反的那件事摆正。

\begin{center}
\begin{tikzpicture}[x=1cm,y=1cm,font=\small,>=stealth]
  \draw[black!60] (6,-0.15) -- (6,5.95);
  \node[font=\footnotesize,black!80,anchor=south] at (6,6.0) {真值（未知且固定）};
  \foreach \y/\a/\b in {0.25/5.05/7.05, 0.57/4.60/6.60, 0.89/5.30/7.30,
    1.21/4.90/6.90, 1.53/5.50/7.40, 1.85/4.35/6.45, 2.49/5.15/7.15,
    2.81/4.75/6.85, 3.13/5.40/7.20, 3.45/4.50/6.50, 3.77/5.20/7.10,
    4.09/4.95/6.95, 4.41/5.60/7.35, 4.73/4.40/6.40, 5.05/5.35/7.25,
    5.37/4.85/6.75, 5.69/5.10/7.00}
    {
      \draw[blue!60!black] (\a,\y) -- (\b,\y);
      \draw[blue!60!black] (\a,\y-0.07) -- (\a,\y+0.07);
      \draw[blue!60!black] (\b,\y-0.07) -- (\b,\y+0.07);
    }
  \draw[red!75!black,very thick] (6.25,2.17) -- (8.15,2.17);
  \draw[red!75!black,very thick] (6.25,2.09) -- (6.25,2.25);
  \draw[red!75!black,very thick] (8.15,2.09) -- (8.15,2.25);
  \node[align=left,anchor=west,font=\footnotesize,black!75] at (8.7,4.1)
    {每一条是重复一次实验\\造出来的区间};
  \node[align=left,anchor=west,font=\footnotesize,red!70!black] at (8.7,2.17)
    {这一次没盖住——\\但事后无从分辨};
  \node[align=left,anchor=west,font=\footnotesize,black!75] at (8.7,0.6)
    {\(95\%\) 说的是这一列的\\长期盖住比例};
\end{tikzpicture}

\emph{真值是那条不动的竖线，随机的是每一条横线。\(95\%\) 描述的是这一整列区间的行为，而不是其中任何一条与竖线的关系。}
\end{center}

\hyperref[sec:11-Mathematical-Statistics/04-Interval-Estimation/01]{``置信度属于程序，而非这次参数''}给出的判断就是这张图：\(95\%\) 说的是造区间这套程序在重复使用下有 \(95\%\) 的时候会盖住真值；对手上这一个区间，真值要么在里面要么不在，没有中间状态。想说``参数落在区间里的概率是 \(0.95\)''，那是可信区间的语言，需要先给出先验，得到的也是另一个东西。

同样的方向问题出现在 \(p\) 值上。\hyperref[sec:11-Mathematical-Statistics/05-Hypothesis-Testing/02]{``\(p\) 值不是原假设为真的概率''}把它写成两个条件概率的区别：\(p\) 值是 \(P(\text{数据这么极端}\given H_0)\)，而人们想知道的是 \(P(H_0\given\text{数据})\)。两者由 Bayes 公式相连，中间隔着一个先验；先验很低时，\(p=0.04\) 对应的后验可以仍然偏向原假设。忽略这一点，会让每一次``显著''都被当成``证实''。

这两条判断能站住，都是因为背后有一个思想实验：\hyperref[sec:11-Mathematical-Statistics/01-Statistical-Models/03]{``抽样分布连接数据与推断''}讲的``如果重复整个实验，这个统计量会怎样分布''。我们只有一份数据，抽样分布描述的是那些没有发生的平行世界——而正是它把一次观测变成了一句带范围的断言。上图那一列横线画的就是这些平行世界。

\subsection{被评估的是整条流程，不是那个模型}

第二件容易搞混的事是评估对象。交叉验证给出的数字不属于模型对象，属于从原始数据到预测的那条完整流程。\hyperref[sec:13-Machine-Learning/01-Learning-Problems-and-Evaluation/04]{``数据划分、泄漏与分布偏移''}把这条讲成一个可操作的要求：凡是要用到标签或全量统计量的步骤——标准化、缺失填充、特征选择、目标编码、超参搜索——都必须搬进折内重做。任何一步站在折外，得到的就是一个系统性乐观的数，而且乐观的幅度无法事后估计。

上一篇末尾提到的种子问题在这里有了对应：如果两个方案的差距小于同一方案在不同划分下的波动，那么这个比较不含信息。\hyperref[sec:11-Mathematical-Statistics/06-Resampling-and-Model-Selection/01]{``Bootstrap 近似估计量的抽样波动''}提供的做法是让数据自己给出这个波动幅度——重复重采样、重跑整条流程、看指标分布有多宽。这个数应当和方案之间的差距放在一起看，单独报一个点估计等于隐去了比较的分母。

\hyperref[sec:11-Mathematical-Statistics/06-Resampling-and-Model-Selection/02]{``置换检验依赖可交换性''}给出的是另一件同样便宜的工具，用在评估上就是那个标签打乱自检：把标签随机打乱，整条流程原样重跑，指标应当回到随机水平。回不去，说明流程里存在让标签信息绕过模型的通道。这个检查抓不到所有泄漏，但它抓到的那些，通常就是离线线上脱节最严重的那几种。

\subsection{一次检验的保证，在反复使用之后就不再是那个保证}

回到开头。名义错误率是对单次判定的承诺，而实际工作中很少只判定一次：同时试十五个特征组合、比较八个模型、看四个细分人群、每天窥视一次实验结果。\hyperref[sec:11-Mathematical-Statistics/05-Hypothesis-Testing/04]{``多重比较改变错误的计量单位''}给出的判断是，这时该控制的量已经不是单次错误率，而是整族判断里的错误比例——不改口径就报显著，等于把``二十次里错一次''当成``每次都对''。边看边停是同一个漏洞的连续版本：给足够多的观察时机，游走的统计量迟早会碰到阈值。

功效是另一半，通常被完全省略。\hyperref[sec:11-Mathematical-Statistics/05-Hypothesis-Testing/01]{``两类错误与功效共同定义检验''}的要点是只控制第一类错误等于只承诺``不冤枉好人''，不承诺``抓得住坏人''。这带来一个反直觉但很重要的后果：低功效的实验不仅容易漏掉真效应，它\emph{通过}显著性筛子的那些结果还会系统性地夸大效应量——因为在样本量不足时，只有恰好抽到偏大的一次波动才可能越过阈值。所以``小样本上做出来的显著提升''往往在放大样本后缩水，这不是运气变差，是筛选机制的必然。开头那个案例里 \(2.1\%\) 的离线提升上线后归零，多半就是这个形态叠加了边看边停。

\subsection{泛化界告诉你什么，不告诉你什么}

理论这一侧提供的是另一类保证。\hyperref[sec:13-Machine-Learning/10-Learning-Theory/01]{``训练误差为何总是乐观''}给出了根本原因：参数是按训练集挑的，于是训练误差不再是真风险的无偏估计，偏低的幅度随可选假设的多少而增长。这句话本身就解释了为什么必须有一个从未参与选择的留出集。

至于那些形如``真风险不超过经验风险加上一个复杂度项''的上界，\hyperref[sec:13-Machine-Learning/10-Learning-Theory/04]{``泛化界的解释力与边界''}说得很诚实：代进真实数字后，界常常给出大于 \(1\) 的误差上限，也就是毫无内容。它们的价值不在数值，而在于指出哪些量控制泛化——样本量、假设类容量、间隔、稳定性。把界当成预测工具会失望，把它当成``该往哪个方向使劲''的指南则相当可靠。

需要有限样本、可用数值的保证时，\hyperref[sec:11-Mathematical-Statistics/07-Statistical-Decision-and-Reliable-Prediction/05]{``Conformal prediction 用可交换性换有限样本覆盖''}是少数几个真的能拿来用的工具：它不假设分布形式，只要求校准数据与待预测样本可交换，就给出有覆盖保证的预测集合。代价写在那个前提里——可交换性一旦被时间漂移破坏，保证同时失效，而这正是下一篇要处理的事。

\subsection{这一步做错时，故障长什么样}

A/B 实验反复不显著又反复上线，说明决策规则实际上是``看到好数就停''，而不是事先定好的那条。这种流程长期看会让一批实际无效的改动累积进系统，而每一个改动当时都有``显著''的证据。

离线提升在线上消失，前面已经出现过两次，但它在这一步有第三种成因：离线评测集与线上流量的分布不同，或者离线指标与线上指标度量的根本不是同一件事（离线看排序质量，线上看留存）。区分它与泄漏的办法是看差距是否稳定：泄漏造成的差距从第一天起就固定在那里，分布不同造成的差距会随时间和人群变化。

同一个模型换个评测集排名反转，说明评测集之间的差异大于模型之间的差异。此时正确的结论不是``某个评测集不准''，而是这两个模型在当前证据下无法区分。硬要选一个，选失效方式已知的那个。

指标全都很好，但没有人能说清这些数是在什么条件下成立的——这是最危险的形态，因为它不产生任何症状，直到某天条件变了。评估报告里真正该被审视的不是那几个数，是它们下面那行常常被省略的``在什么条件下''。

到这里为止，整条链路还停留在一个预测问题上：世界在那里，系统去猜它。上线改变了这个前提——系统的输出开始进入世界，然后回到它自己的输入。
